\chapter{Sistemas de Control Estructural}

% NOTA: Utilizar apuntes de Manuel Ruiz Sandoval como guía
% Papers: 
% - Symans et al. 2008 - Energy Dissipation Systems for Seismic Applications: Current Practice and Recent Developments
% - Housner et al, 1997 - Structural Control. Past, Present, and Future
% - Saaed et al, 2015 - A State-of-the-art review of structural control systems
% - 

%% Anotaciones - Ideas
% En sistemas de control estructural, en introducción, poner la formula general de un sistema con control Mxpp + Cxp + Kx = F_externas + F_control
% Explicar que F_control tiene varias formas de diseñarse
% En vez de explicar los dispositivos de control en secciones de control pasivo, activo, semi-activo e híbrido... poner una sección de tipos de control explicandolos brevemente.
% Luego, crear secciones de modelos de sistemas de control, explicar brevemente como se modelan y ecuaciones de TMD, Amortiguador viscoso, Disipador histerético, BRB, 
%

\section{Introducción}
% Some from Saaed et al 2015
Los sistemas de control estructural son dispositivos mecánicos que se incorporan a una estructura para poder mejorar su desempeño y serviciabilidad ante cargas dinámicas.

\begin{itemize}
    \item Reducir la \underline{respuesta}
    \item Reducir \underline{pérdidas por daños}
    \item Reducción de \underline{riesgos} (Ej. Probabilidad de colapso)
\end{itemize}

\vspace{1eX}

Se clasifican en los siguientes grupos:

\begin{description} 
    \item[Pasivos] Estos sistemas no requieren una fuente de energía externa para operar. Funcionan disipando o absorbiendo la energía que ingresa a la estructura basándose únicamente en sus propiedades físicas y en el propio movimiento de la estructura.
    \item[Activos] Estos sistemas requieren una fuente de poder externa significativa (electricidad, sistema hidráulico) para funcionar. Utilizan sensores que miden la respuesta de la estructura en tiempo real y un algoritmo de control que calcula una fuerza de corrección. Luego, los actuadores aplican activamente esa fuerza a la estructura para contrarrestar la vibración.
    \item[Semi-Activos] Estos sistemas representan un punto intermedio inteligente. Requieren una pequeña cantidad de energía (como la de una batería), pero esta energía no se usa para aplicar fuerzas directamente a la estructura (como en el control activo). En cambio, se utiliza para cambiar las propiedades mecánicas (rigidez o amortiguación) del dispositivo de control en tiempo real.
    \item[Híbridos] Como su nombre indica, estos sistemas combinan dos o más de los tipos anteriores para aprovechar sus ventajas.
\end{description}


En general, las estructuras con dispositivos de control tienen la siguiente EDM:

\begin{align*}
    {M}{\ddot{q}} + {C}{\dot{q}} + {K}{q} + {f}_{p} + f_{s} = {f}_a + {f}_e
\end{align*}

\noindent donde ${f}_{p}$ es la fuerza de dispositivos pasivos, ${f}_{sa}$ es la fuerza de dispositivos semi-activos, ${f}_a$ es la fuerza de dispositivos activos, y ${f}_e$ es la fuerza externa (p.ej., terremoto).

\section{Control pasivo}
Los sistemas de control pasivo no requieren de una fuente de energía externa. \\

Funcionan absorbiendo, aislando y disipando energía. \\

Su operación puede depender tanto de el desplazamiente, velocidad y aceleración relativa entre los nodos de conexión del dispositivo. Algunos de estos dispositivos requieren ser reemplazados luego de ser utilizados. \\

% La ecuación general para un sistema estructural con dispositivos de control pasivo es:

% \begin{align}
%     {M}{\ddot{q}} + {C}{\dot{q}} + {K}{q} + \textcolor{red}{{f}_p(q,\dot{q},\ddot{q})} = {f}_e
% \end{align}

\subsubsection{Amortiguadores viscosos}

\begin{equation}
    f_p(\dot{r}) = c_d |\dot{r}|^\alpha \text{sign}(\dot{r})
\end{equation}

\subsubsection{Amortiguadores metálicos}

Disipación de energía por histeresis. Por ejemplo, usando un modelo de Bouc-Wen:

\begin{align}
    f_P(r,\dot{r},t) = (1-\alpha) k r +  \alpha k z(t)
\end{align}

donde $z(t)$ es un estado interno que maneja la histeresis del dispositico, y que evoluciona a través de una ecuación diferencial nolineal:

\begin{align}
    \dot{z} = h(\dot{r},z)
\end{align}

\subsubsection{Amortiguadores de fricción}

\begin{align}
    f_P(\dot{r},t) = \mu(t) N(t) \text{sign}(\dot{r})
\end{align}

donde $\mu(t)$ es el coeficiente de fricción dinámico, $N(t)$ es la fuerza normal en la superficie de deslizamiento, y $\dot{r}$ es la velocidad relativa entre los terminales del amorgtiguador de roce.

% \subsubsection{Dispositivos Viscoelásticos}

% Comunmente son pistones que consisten en cilindros huecos rellenos de un material viscoelástico (fluido o sólido). En el caso de meteriales fluidos, la fuerza se genera cuando el fluido orificios y en el caso de materiales sólidos ocurre cuando estos se deforman. 

% Para modelar las fuerzas:
% \begin{align*}
%     P(t) &= C|\dot{u}(t)|^\alpha \text{sign}[\dot{u}(t)] \longrightarrow \text{VE Fluidos} \\
%     P(t) &= Ku(t) + C\dot{u}(t) \longrightarrow \text{VE Sólidos}
% \end{align*}
% La disipación de energía ocurre por la fricción generando calor.


\subsection{Amortiguadores de masa sintonizada}

\begin{align}
    f_P = m_{TMD} a_{TMD}
\end{align}

donde $a_{TMD}$ es la aceleración absoluta del TMD.

% \subsubsection{Otros}

% \begin{itemize}
%     \item Dispositivos de re-centrado
%     \item Amortiguadores de transformación de fase
%     \item Amortiguadores/disipadores de vibración dinámica:
%     \begin{itemize}
%         \item Amortiguadores de Masa Sintonizada (TMD)
%         \item Amortiguadores de Líquido Sintonizado (TLD)
%     \end{itemize}
% \end{itemize}


% \subsubsection{Dispositivos de Aislación Sísmica}
% % Saaed

% \begin{itemize}
%     \item Eficiente para vibraciones desde el suelo como sismos y tráfico, no para viento.
%     \item Alta flexibilidad lateral
%     \item Alta rigidez vertical
%     \item Inclusión de nuevo modo de vibrar con \emph{Desplazamientos relativos entre piso} pequeños.
%     \item Al incluir el nuevo modo de vibrar, se alejan los otros modos lejos del contenido de periodos de la excitación.
%     \item Funcionan para edificios bajos con periodos fundamentales cortos y no para edificios altos que tienen periodos fundamentales largos.
%     \item Se puede incluir en distintas ubicaciones de la estructura.
% \end{itemize}

% Los dispositivos de aislación sísmica basan su funcionamiento en la flexibilidad horizontal de un nivel pero una alta rigidez vertical. Debido a la flexibilidad horizontal, se introduce un nuevo modo de vibrar que contiene \emph{desplazamientos relativos entre pisos} pequeños. 

\section{Control activo}

Los dispositivos de control activo requieren de energía para contrarrestar la respuesta estructural. Funcionan adaptandose en función de la excitación y respuesta. Requieren del uso de sensores y actuadores. La ecuación general es la siguiente:

\begin{align}
    {M}{\ddot{q}} + {C}{\dot{q}} + {K}{q} = \textcolor{red}{{f}_a(q,\dot{q},\ddot{q},u)} + {f}_e
\end{align}

donde $f_a$ es la fuerza activa, que puede depender de estados del sistema, pero depende principalmente de una señal de control exógena $u(t)$.

\underline{Ventajas del control activo}:
\begin{itemize}
    \item Efectividad, no hay límites teoréticos para la eficiencia del controlador
    \item Adaptabilidad a la excitación
    \item Selectvidad, el controlador se puede diseñar para varios objetivos (seguridad y comfort)
    \item Aplicabilidad para un gran rango de frecuencias y excitaciones
\end{itemize}

\underline{Dispositivos}:
\begin{itemize}
    \item Amortiguador de masa activo (Active mass damper, AMD)
    \item Sistemas de tendones activos (Active tendon systems)
    \item (Active brace systems)
    \item (Pulse generation systems)
\end{itemize}

\section{Control semi-activo}

\begin{align}
    {M}{\ddot{q}} + {C}{\dot{q}} + {K}{q} + \textcolor{red}{{f}_s(q,\dot{q},\ddot{q},u)} = {f}_e
\end{align}

Ejemplo: Amortiguador viscoso de orificio variable. El coeficiente $c_d(u)$ es una función de una señal de control exógeno.

\begin{equation}
    f_{sa}(\dot{r},u) = c_d(u) |\dot{r}|^\alpha \text{sign}(\dot{r})
\end{equation}

\begin{itemize}
    \item TMD semi-activo (Semi-active TMDs)
    \item TLD semi-activo (Semi-active TLDs)
    \item Amortiguador de fricción semi-activo (Semi-active friction damper)
    \item (Semi-active vibration absorbers)
    \item Controlador de rigidez semi-activo
    \item Amortiguadores Electroreologicos (Electrorheological damper, ER)
    \item Amortiguador Magnetoreologico (MR)
    \item Amortiguador viscoelástico semi-activo
\end{itemize}

\section{Control híbrido}

\begin{itemize}
    \item Amortiguador de masa hibrido (Hybrid mass damper, HMD)
    \item Aislación basal híbrida: aislación basal + control activo o semi-activo
\end{itemize}
